% =====================================================================
%  Consentimiento informado -- persona evaluadora del cuasi-experimento
%  Proyecto SIMPA -- Equipo AHMRV -- ISR-401 -- UTEQ -- 2026-2027 PPA
%  Compilacion: pdflatex consentimiento_evaluador.tex  (x2)
% =====================================================================
\documentclass[11pt,a4paper]{article}

\usepackage[utf8]{inputenc}
\usepackage[T1]{fontenc}
\usepackage[margin=2.2cm]{geometry}
\usepackage{array}
\usepackage{booktabs}
\usepackage[table]{xcolor}
\usepackage{enumitem}
\usepackage{tcolorbox}
\usepackage{amssymb}
\usepackage{titlesec}

\definecolor{uteqverde}{HTML}{1F6F3C}
\definecolor{grisclaro}{HTML}{F2F2F2}

\newcolumntype{L}[1]{>{\raggedright\arraybackslash}p{#1}}
\newcolumntype{C}[1]{>{\centering\arraybackslash}p{#1}}

\titleformat{\section}{\normalfont\normalsize\bfseries\color{uteqverde}}{\thesection.}{0.5em}{}
\titlespacing*{\section}{0pt}{9pt}{3pt}

\setlength{\parskip}{3pt}
\setlength{\parindent}{0pt}
\pagestyle{empty}

\newcommand{\casilla}{$\square$}

\begin{document}

\begin{center}
{\footnotesize UNIVERSIDAD TÉCNICA ESTATAL DE QUEVEDO\\
Facultad de Ciencias de la Computación --- Carrera de Software}\\[8pt]
{\large\bfseries\color{uteqverde} CONSENTIMIENTO INFORMADO}\\[2pt]
{\normalsize Persona evaluadora --- estudio de calidad de requisitos}\\[2pt]
{\footnotesize Protocolo registrado en OSF \texttt{4z35d}}
\end{center}

\vspace{2pt}

\section{En qué consiste su participación}

Se le pide leer un conjunto de \textbf{50 requisitos funcionales} de un sistema
de gestión de cultivo de palma africana y puntuarlos en cinco dimensiones de
calidad, siguiendo una rúbrica que se le entrega y se le explica antes de
comenzar.

Los requisitos provienen de dos fuentes distintas y \textbf{usted no sabrá cuál
corresponde a cada uno}. Esta condición es parte del diseño del estudio: su
finalidad es que la valoración no se vea influida por la procedencia. Al
terminar la sesión se le informará de cuáles eran esas fuentes.

Duración estimada: entre 60 y 90 minutos, con los descansos que necesite.

\section{Qué se hace con sus respuestas}

\begin{itemize}[nosep,leftmargin=1.5em]
  \item Sus puntuaciones se identifican con un \textbf{código} (E1, E2, E3\dots),
        nunca con su nombre.
  \item Se analizan de forma agregada junto con las de las demás personas
        evaluadoras, y se comparan entre sí para medir el grado de acuerdo.
  \item Los datos anonimizados pueden depositarse en un repositorio público de
        datos abiertos como parte del paquete de replicación del proyecto.
  \item Su nombre no aparece en ningún archivo publicado.
\end{itemize}

\section{Qué NO se evalúa}

\begin{tcolorbox}[colback=grisclaro,colframe=uteqverde,boxrule=0.6pt,arc=2pt]
\textbf{No se le evalúa a usted.} Sus puntuaciones no son un examen, no se
califican, no se comunican al profesorado de ninguna asignatura y no tienen
efecto alguno sobre su situación académica. Lo que se estudia es la calidad de
los requisitos, no su criterio.
\end{tcolorbox}

\section{Carácter voluntario}

Su participación es voluntaria. Puede interrumpirla en cualquier momento y sin
dar explicaciones, y puede pedir que sus datos se retiren del estudio hasta el
momento en que se publiquen de forma agregada.

No recibe compensación económica ni beneficio académico por participar.

\section{Riesgos}

No se prevén riesgos derivados de esta participación. La tarea consiste en leer
y puntuar textos técnicos, y no se recogen datos personales sensibles.

\section{Autorizaciones}

Marque lo que autorice. \textbf{Las tres son independientes} y puede aceptar
unas y rechazar otras.

\vspace{3pt}
\renewcommand{\arraystretch}{1.5}
\begin{tabular}{C{1cm} L{15cm}}
\casilla & Autorizo el uso de mis puntuaciones, de forma anonimizada, en el
análisis del estudio. \\
\casilla & Autorizo que mis puntuaciones anonimizadas se publiquen como parte
del paquete de datos abiertos del proyecto. \\
\casilla & Autorizo que se conserve mi hoja de puntuación firmada como evidencia
del estudio, en zona restringida y cifrada. \\
\end{tabular}

\vspace{4pt}
{\footnotesize Si no marca la primera casilla, sus puntuaciones no se incorporan
al análisis y la hoja se destruye al terminar la sesión.}

\section{Tratamiento de datos personales}

El único dato personal que se recoge es su nombre y su firma en este documento,
con la finalidad exclusiva de acreditar el consentimiento. Ese documento se
conserva en zona restringida, cifrado, y no se publica. Sus puntuaciones se
disocian de su identidad mediante un código antes de cualquier análisis.

Puede ejercer sus derechos de acceso, rectificación y supresión dirigiéndose al
equipo responsable por el correo institucional que consta al pie.

\section{Declaración}

He leído este documento, he podido preguntar lo que he querido y he recibido
respuesta. Acepto participar en los términos que he marcado.

\vspace{6pt}
\renewcommand{\arraystretch}{1.4}
\begin{tabular}{L{5.5cm} L{5.2cm} L{5cm}}
\toprule
\rowcolor{grisclaro}
\textbf{Nombre completo} & \textbf{Código asignado} & \textbf{Fecha} \\
\midrule
\rule{0pt}{4.5ex} & & \\
\bottomrule
\end{tabular}

\vspace{20pt}

\begin{center}
\begin{tabular}{C{7cm} C{7cm}}
\rule{6.2cm}{0.4pt} & \rule{6.2cm}{0.4pt} \\
{\footnotesize Firma de la persona evaluadora} &
{\footnotesize Firma de quien modera la sesión} \\
\end{tabular}
\end{center}

\vfill
{\footnotesize
Equipo AHMRV --- Ingeniería de Requerimientos (ISR-401) --- Universidad Técnica
Estatal de Quevedo --- Período 2026--2027 PPA.\\
Docente responsable: Ing. Gleiston Cicerón Guerrero Ulloa, PhD ---
\texttt{gguerrero@uteq.edu.ec}}

\end{document}
